\documentclass[11pt,oneside]{article}
\pdfoutput=1
\usepackage[utf8]{inputenc}
\usepackage[T1]{fontenc}
\usepackage{lmodern}
\usepackage[a4paper,margin=1.05in]{geometry}
\usepackage{amsmath,amssymb,amsthm}
\usepackage{microtype}
\usepackage{enumitem}
\usepackage{booktabs}
\usepackage{array}
\usepackage{longtable}
\usepackage{xcolor}
\usepackage{titlesec}
\usepackage{fancyhdr}
\usepackage{setspace}
\usepackage[hidelinks,breaklinks]{hyperref}
\usepackage{parskip}
\definecolor{accent}{RGB}{40,72,120}
\hypersetup{colorlinks=true,linkcolor=accent,citecolor=accent,urlcolor=accent}
\titleformat{\section}{\large\bfseries\color{accent}}{\thesection}{0.6em}{}
\titleformat{\subsection}{\normalsize\bfseries}{\thesubsection}{0.6em}{}
\pagestyle{fancy}
\fancyhf{}
\fancyhead[L]{\small\itshape What Computation Costs}
\fancyhead[R]{\small\thepage}
\renewcommand{\headrulewidth}{0.4pt}
\theoremstyle{definition}
\newtheorem{test}{Test}
\theoremstyle{remark}
\newtheorem{pred}{Prediction}
\newtheorem{hyp}{Hypothesis}
\setstretch{1.0}
\setlength{\parskip}{0.4em}
\newcommand{\qt}[1]{\begin{quote}\small #1\end{quote}}
\title{\vspace{-1.2cm}\textbf{\color{accent}What Computation Costs}\\[0.4cm]
\large Software engineering as the work of an interpreter\\
standing between a person and a Turing machine}
\author{
Edwin O Onyango\\[-1pt]
\small\texttt{edwin.o.onyango.jr@gmail.com}
\and
Walter O Wagunde\\[-1pt]
\small\texttt{wagudewalter2@gmail.com}
}
\date{September 2026}
\begin{document}
\maketitle
\thispagestyle{empty}
\vspace{-0.5cm}
\begin{abstract}
\noindent In 1936 Turing described a computer by watching a person do arithmetic and
setting every cost that person pays to zero. Software engineering then grew up between the
machine and the people who want things from it, and this paper describes software
engineering from the outside. Code is the layer through which people tell the machine what
to do. A software engineer is an interpreter, standing between a person who wants something and
a machine that does exactly what its sheet says. The interpreter's work has three tiers:
finding out what is wanted, turning it into a \emph{software intention}, and turning the
intention into code. Code generators are absorbing the third tier. Software engineering keeps most of
what it knows in the second tier. It keeps it as rules, and the conditions under
which each rule holds were dropped as the rule was passed on, which is why every
experienced engineer eventually learns that it always depends.

The claim of the paper is that those dropped conditions can be recovered. They were never
anything but the machine's constraints, and the machine's constraints come in two kinds.
A few things it cannot do at any price; those are walls, and all but one are theorems. Everything
else it can do at a price, and the prices are exactly the eight costs Turing idealised
away. The eight sort into three groups. Five are physical, and physics sets floors under them. One is
the boundary between the machine and the people who write for it: somebody has to write
the program. Two are the world: the input comes from people, and the job never ends. The
field's cost models stop after the physical five, its theorems stop after the boundary one,
and beyond that line there are only rules.
In Turing's machine every move is completely determined by the state and the symbol under
the head; nothing is ever chosen, because nothing costs anything, and a choice exists only
where two options cost different amounts. So every engineering decision is a choice between
alternatives that pay differently under one or more of the eight. That gives every decision
a shape, and every rule a test: what is it answerable to? A rule can answer to a cost, to
a theorem, to a measurement, or to nothing but the fact that others follow it. Put to four
rules everyone knows, the test finds one resting on a formula that has been recomputed for
forty years, one resting on a cost its own author named and nobody quotes, one resting on
a theorem that cannot license it, and one whose number is the width of a 1928 punched
card, kept because everyone keeps it. A decision
itself has a cost with four parts: writing the change, checking it, the chance it is wrong
times the cost of undoing it, and what is spent finding out before committing. Code
generators remove the first part. The three that remain are all about being wrong, so what
has to be verified is the decision, not the code. The paper ends with seven predictions and an invitation to
find a principle that does not fit.
\end{abstract}
\vspace{0.2cm}
\hrule
\vspace{0.4cm}

% ===========================================================================
\section{The Machine}\label{sec:machine}
% ===========================================================================
In 1936 a ``computer'' was a person at a desk with paper, a pencil and a printed sheet of
instructions, doing work that astronomers, insurers and the makers of artillery tables
depended on. Turing was answering Hilbert's question of whether a mechanical procedure
could settle any mathematical statement, and had to make ``mechanical'' precise. He did it
by describing what that person does \cite{turing1936,hodges1983}, and the description was
exact enough that he later defined a digital computer against it \cite{turing1950}. There
is a tape of squares, a head over one square, a finite state, and a sheet that says, for
each state and symbol, what to write, which way to move, and what state to enter.

Is that all a person computing by rule ever does? Turing claimed so. The claim cannot be
proved, only searched for exceptions, and the search has failed. Every other definition of
computation offered before or since computes the same class; Church's calculus and Gödel's
general recursive functions, which Kleene showed equivalent, were among the first
\cite{church1936,kleene1936}. That is why it is the Church-Turing \emph{thesis}. Of those
definitions, Turing's is the one that describes a person: the lambda calculus has no tape
and no head, hence no notion of where a value is or what reaching it costs, and that is
the feature that made Turing's analysis persuasive where Church's was not
\cite{gandy1980,sieg1997,soare1996}. The paper uses the thesis twice. Read, write, move,
change state and halt are taken as the whole of what any machine does; those are the five
operations in the model's modern statement, where halting is a state, and Turing's own
machines printed for ever. And a claim held by a failed search for exceptions is the
standard the paper's own lists are held to.

Every operation costs the clerk something, and the desk imposes five costs the model
omits: paper must be bought, a distant square takes longer to reach than the next one,
one person sits at the desk, paper smudges and memory fails, and anyone who walks past can
read what is on it. Three more things cost him and are free in the model. Somebody wrote
the sheet. Somebody decided what goes on the tape. And the work is never done once, since
a table is corrected for as long as anyone uses it, whereas in the model the job is fixed
when the machine starts. Turing set all of it to zero because what is computable cannot
depend on how fast a pencil moves.

Philosophy of science has a name for this move. An \emph{idealisation} is a real quantity
assumed to be zero so that a theory becomes tractable: the frictionless plane, the point
mass, the ideal gas. Putting the quantity back afterwards, to see what it changes, is a
method in its own right \cite{mcmullin1985}. Turing's machine is the ideal gas of
computing, and if the costs of computing are to be put back into one model, it has to be
this one. To idealise a cost away Turing first had to assume something in its place, one
tape, one head, a sheet already written, a job that is fixed, so the model is an inventory
of assumptions, each standing in for a cost the clerk pays, and an inventory can be
audited one line at a time. Dally \cite{dally2022}, Valiant \cite{valiant1990} and the
Fallacies of Distributed Computing \cite{fallacies} each restored the costs they needed and
stopped. This paper reads the audit to the end, and says what a complete audit is for.

One more feature of the model carries \S\ref{sec:decisions}. In an automatic machine, in
Turing's words, ``the motion of a machine is completely determined by the configuration'':
the state and the symbol under the head fix the next move, and there is no point at which
anything is chosen. Turing did allow a machine that at some configurations ``cannot go on
until some arbitrary choice has been made by an external operator'', and set it aside at
once, ``in this paper I deal only with automatic machines'' \cite{turing1936}; and the
programmer, who chose what to write on the sheet, is outside the model altogether. So the
model removes choice twice, and every decision a real engineer makes entered from one of
two places: the writing of the sheet, or a cost the model set to zero. The real clerks were
already paying under both. Prony decided how the sheet was written when he split
arithmetic into steps unskilled workers could follow, and the remedy for errors in printed
tables, two computers working the same sheet, is redundancy against a tape that fails
\cite{grier2005}. Engineers decide the same things at speeds eight orders of magnitude
higher, plus the one thing the clerks never decided: the sheet. How anyone writes a sheet
is the next question.

% ===========================================================================
\section{Code}\label{sec:code}
% ===========================================================================
Nobody writes Turing's sheet. Code is a stack of translations that ends at it: framework,
library, language, compiler, instruction set, circuit. Every layer exists for one reason,
so that a person can say more per symbol written and read more per symbol read. The
machine can already do anything any layer expresses. So the layers change nothing about
what can be said; they change only what it costs a person to say it and to understand it.
That cost is one of the eight the model removed; idealisation~6, below, is the assumption
that hides it.

At every layer the meaning of a text is fixed. A language standard fixes what a text means
by fixing what the machine must observably do, in the C standard's words ``the observable
behavior of the abstract machine'' \cite{c11}. This is a claim about meaning, not about
execution: a program may draw random numbers or run threads in any order, and that is
meaning fixed as ``any of these'', on purpose. The defect the profession fences off is
meaning left unfixed by accident, undefined behaviour, a race nobody intended. A learned model is a program in this sense too, fixed weights with fixed meaning on the
same hardware, and what varies in its output is sampled on purpose. A text with one
meaning can be checked without being run, and every practice of reading, reviewing, typing
and proving code rests on that.

So code is a way of telling the machine what to do, in a form a person can write, that
does exactly what it says. The computer collapsed the price of executing a procedure once
it was specified; what code cannot supply is the specifying, what to say. The people who
want something from the machine do not speak in sheets, and between them and the machine
a software engineer has to stand.

% ===========================================================================
\section{The Interpreter}\label{sec:interpreter}
% ===========================================================================
A software engineer is an interpreter, in the courtroom sense rather than the compiler
sense. An interpreter renders meaning live between two parties who share no language,
chooses words under the constraints of the second language, and is answerable to both
sides for the result. On one side here is someone who wants something and can say so only
in the terms of their own world: an invoice must not be altered after it is sent, a
customer should not wait. On the other side is a machine that will do exactly what its
sheet says, at the prices the world charges for each line, and that understands nothing
else. Everything software engineering knows is knowledge about how to stand between them.

Before any code is written the engineer produces a statement of what the code must do,
and the paper calls it a \emph{software intention}: a statement of what the machine is to
be told, in terms the machine can be told, made before the telling. It is not the want,
which lives on the world side. It is not the code, which is one of many texts that could
carry it. It is where the decisions live: which quantity must be preserved and at which
level, what may be assumed about the world, what the machine will remember and for how
long, where the boundaries fall, what happens when a part fails. The idea has neighbours.
Simon described any artifact as a meeting point between an inner environment, its own
substance, and an outer environment, the surroundings it must serve, ``an interface in
today's terms'' \cite{simon1996}; this paper is about the engineer who builds the meeting
point, and what the inner environment charges. Brooks's essence is the conceptual
construct that is left when the accidents of representation are stripped away, and he
predicted that only the accidents would shrink \cite{brooks1987}. Jackson and Zave put a
specification at the interface between world and machine, with the requirements on the
world side \cite{zavejackson1997}; the reference model built on that, world, requirements,
specification, program, machine, is the layering the tiers below inherit \cite{gunter2000}.
Naur's theory is what the builders hold and the text does not \cite{naur1985}. The new
name claims that the software intention is where the engineer's expertise is, and that it
has coordinates, which \S\ref{sec:decisions} gives.

\subsection{Three Tiers}\label{sec:tiers}

The work has three tiers, and each is a different translation. The first tier goes from
the world to the requirements: finding out what is wanted, what the people who want the
software actually want, which constraints are real, what will change. The second tier goes
from the requirements to the specification, and on to the design decisions the reference
model leaves inside the program. It turns what is wanted into a software intention, and it
is a negotiation with the machine, which can do a few things at no price whatever and
everything else at a price. The third tier goes from there to the program text. It is the
tier software engineering has spent most of its history on, and the tier code generators
are absorbing. The layering is the reference model's. This paper says what shapes the
second tier, and what happens to the third when it becomes free. The paper leaves the
first tier alone; \S\ref{sec:idealisations} says why.

The lines between the tiers move: the line between the second and the third is the cost
of producing code, and \S\ref{sec:formula} is about its collapse. What does not move is that
the middle tier is where the decisions are, and that every decision there is made against
the machine's constraints. So the paper needs two things: a sorting of what the
interpreter already knows, and a complete list of the constraints. The first points at
the second.

% ===========================================================================
\section{Three Kinds of Knowledge}\label{sec:knowledge}
% ===========================================================================
Knowledge of software comes in three kinds, told apart by what happens when a piece is
handed to somebody else. \emph{Theorems} regenerate: whoever holds the reason can rebuild
every consequence anywhere, because the conditions travel inside the statement.
\emph{Lookups} do not regenerate at all: a port number, a flag, the name of a method
somebody else wrote. \emph{Rules} sit between: index the column the query filters on, keep
services small, pin the dependencies, do not optimise prematurely. A rule generates as a
theorem does, one sentence covering a thousand situations, but it generates correctly only
under conditions, and the conditions are not in the sentence. Whoever found the rule found
it as the cheapest answer under a particular machine, team, business and set of prices. The
answer was passed on and the conditions stayed behind, usually because they were too
obvious at the time to state. The three form a spectrum, from statements that carry all
their conditions to statements that carry none, and medicine has the same spectrum, a
mechanism as theorem, a dosage as lookup, a clinical guideline as rule.

The spectrum is not an accident of the field. Knowledge accumulates where reality
answers, and its form is set by how fast and how exactly reality answers. A program holds
still, so a question about programs gets an answer that does not move, sometimes an exact
``no'', and that is a theorem. A fact that holds still but carries no reason, a port
number, regenerates nothing and needs no theory; that is a lookup. A price moves, so a
question about prices gets an answer that is true for a while, and that is a rule whose
conditions are the prices at the time. What people want cannot be measured, so questions
about it get no formal answer at all. Sections \ref{sec:idealisations} and \ref{sec:stops}
check this against the record.

The middle of the spectrum is where the interpreter spends the week, and applying a rule
there is expertise, not recall: the expert supplies the missing conditions from experience,
and cannot always say why. Knuth's remark on premature optimisation is the clearest case
(\S\ref{sec:knuth}): the half that became a proverb has travelled for fifty years, and the
cost it was weighing, two sentences earlier in the same paragraph, has not travelled at
all.

The three must cover what an engineer knows, or what follows is partial, so the
apparent exceptions get a sentence each. Taste is a rule whose conditions have been
compiled beyond the holder's ability to state them (\S\ref{sec:whykeydropped}). A war story
is a rule with a sample of one, and the rare case where the conditions travelled. A
benchmark number is a lookup with a date. A design pattern is a rule published with its
conditions attached, the one place the field has used that form for fifty years
\cite{alexander1977,gamma1994}. A contract is a lookup other people have agreed to keep
true.

\subsection{Method}\label{sec:method}

The paper rests on one observation. A rule is the cheapest answer under some conditions,
and the interpreter's conditions are the machine's constraints, because only the
machine's charges can be written down, and the team and the business enter its account
as the costs of writing the sheet and of the world moving. So the conditions that fell off
every rule were statements about what the machine cannot do and what it charges, and once
the complete list of constraints is recovered, the dropped conditions have somewhere to be
looked up. The list must come from underneath the practice, since a list drawn from the
practice is only the practice again; without it the exercise collapses into ``engineering
involves trade-offs'', which forbids nothing. Software runs on computers, so whatever is
true of every computer is true of every program; and the smallest true description of
every computer is the clerk minus the costs.

The method is de-idealisation, done in full and then used: put each cost back, price it
where physics can, sort the eight, and check the sorting against two counts, one of cost
models and one of theorems; then derive the decision space, put four rules to the test
from their primary sources, and check the result against the field's own history of
practice reversals; then price a decision and remove one term. Claims are marked
\emph{proved}, \emph{measured} or \emph{disputed} where it matters, and every reading of a
source is recorded in a companion document, named at the end. The constraints come first,
and the first of them are the few things the machine cannot do at all.

% ===========================================================================
\section{Walls}\label{sec:cannot}
% ===========================================================================
Walls are few, and all but one are theorems. No machine can decide from another machine's
sheet whether it will run without getting stuck or ever print a given symbol
\cite{turing1936}, and Rice generalised this to every non-trivial question about what a
program computes \cite{rice1953}. A machine with finite memory is formally a finite-state
machine whose halting is decidable by enumerating its states, at a price exponential in
their number, so on a real machine the wall is met as a price nobody can pay, and it
stands. So a sheet cannot in general be checked by reading it, and typing and testing are
the compromises: a type system decides a property conservatively, a test samples
behaviour at points a person chose. Where messages may be delayed without bound, no
deterministic protocol can guarantee to reach agreement if one participant may fail
\cite{flp1985}. A service that must survive a partition cannot, during the partition, be
both consistent and available \cite{gilbertlynch2002}. So every distributed design must
choose what it gives up. And a machine running near the physical limit on its rate of
operation cannot be serial (\S\ref{sec:bills}); that is the one wall that is a physical
limit rather than a theorem.

Two things follow from how short the list is. First, a wall says ``not that'' and never
says which of the remaining options to take, and the remaining options differ only in what
they cost. Even CAP hands the decision to the prices, which decide whether consistency or
availability is given up. Second, walls are the first kind of knowledge, theorems with
their conditions attached, so a practice resting on one can be traced to its fine print.
Everything else the machine can do, at a price, and the prices the universe itself sets
come first.

% ===========================================================================
\section{Four Bills}\label{sec:bills}
% ===========================================================================
A Turing machine costs nothing to define. A physical computer has to be built, powered,
cooled and maintained, and the gap between the two comes apart into four costs, in the
order the model lists its operations: writing, which physics charges for as erasing;
moving; the tape itself, which fails; and the passage of steps. 

\textbf{Erasing.} Forcing an unknown bit to a known value halves the arrangements of the
world consistent with what is known, and the second law makes that entropy appear as heat:
at least $kT\ln 2$ per erased bit, about $2.9\times10^{-21}$ joules at room temperature
\cite{landauer1961}. The cost attaches to erasing, not to computing, since any computation can
run reversibly if a record is kept \cite{bennett1973}. The floor is reached only for
infinitely slow erasure \cite{proesmans2020}. It is \emph{measured}
\cite{berut2012,jun2014,hong2016,yan2018} and \emph{disputed} at its foundations
\cite{earman1999,norton2005}, and nothing here depends on the dispute. Real machines run a
million times above it. This is the famous bill and it bears on nothing anybody does.

\textbf{Moving.} The random access machine cannot be built. It is the model on which every
$O(n\log n)$ is counted, and in it address three costs the same as address three billion. Memory cells have
a minimum size and signals have a maximum speed. So the farthest of $x$ cells packed
around a processor in $d$ dimensions is about $x^{1/d}$ away, the latency to reach the
$x$-th cell is $\Omega(x^{1/d})$, and a pipelined hierarchical memory achieves that bound
\cite{bilardi2009}. This is usually presented as a design heuristic, and it is not one. A
program has a cache because geometry says it must, and no money removes it. Three pieces
of fine print. The bound is on latency to address $x$, not on average latency. The proof
leans on a density bound whose scope is still being worked out \cite{hayden2025}. And it
bounds latency, not throughput, so whether the hierarchy is visible to the programmer
remains a choice; that is where physics ends and the designer's decision begins. 

\textbf{Failing.} More than eight percent of memory modules in Google's fleet saw a
correctable error in a year \cite{schroeder2009}, and flash is comparable \cite{meza2015}.
Below a threshold error rate, majority voting over redundant components buys arbitrary
reliability at finite overhead \cite{vonneumann1956}. The industry pays that overhead as
error correction, replication and scrubbing.

\textbf{Stepping.} A physical system passes through distinguishable states no faster than
a rate set by its energy \cite{margolus1998}. For a kilogram of matter in a litre, the time
for a signal to cross the machine exceeds the time to flip one bit by about $10^{10}$, and
only at black-hole density does the ratio approach one \cite{lloyd2000}. A machine near
that limit must be parallel; that is the physical wall of \S\ref{sec:cannot}. Every machine
ever built runs far below the limit. At network scale the floor is the speed of light. New York to London and back is
near 37 milliseconds in vacuum and 55 in fibre; measured pings are about three times the
vacuum floor, and full fetches about thirty-six times \cite{bozkurt2017}. The gap is
engineering; the floor is not.

\begin{center}
\small
\begin{longtable}{@{}p{0.28\textwidth}p{0.36\textwidth}p{0.28\textwidth}@{}}
\toprule
\textbf{Turing grants for free} & \textbf{The physical bill} & \textbf{Status} \\
\midrule
\endhead
Overwriting a cell is free & $kT\ln 2$ per erased bit, plus a term that grows as erasure is
done faster & Proved; measured; foundations disputed \\
\addlinespace
Any cell is reachable in unit time & $\Omega(x^{1/d})$ latency; a pipelined hierarchy is
optimal & Proved \cite{bilardi2009}; verified here from the abstracts \\
\addlinespace
The tape never fails & Redundancy and correction overhead; $>8\%$ of DIMMs per year &
Measured \cite{schroeder2009}; survivable by \cite{vonneumann1956} \\
\addlinespace
Steps follow one another without limit & A maximum state-change rate per unit energy; the
speed of light; a communication-to-operation ratio near $10^{10}$ & Proved and measured
\cite{margolus1998,lloyd2000} \\
\bottomrule
\end{longtable}
\end{center}



\subsection{Floors and Prices}\label{sec:floors}

A \emph{floor} is a cost nothing lowers, because it follows from the speed of light or the
size of atoms. A \emph{price} is a cost money moves, and prices move by orders of magnitude
within a career. A practice resting on a floor is right for as long as anyone builds
computers. A practice resting on a price has an expiry date, and the value of knowing
which is which is seeing the expiry coming rather than finding it in a postmortem. The
range of prices under which a practice is the cheaper choice is its \emph{regime}.

Physics sets floors and almost never selects a design. Dennard scaling, that shrinking a
transistor's dimensions and voltage together keeps power density constant
\cite{dennard1974}, held until leakage ended it between 2007 and 2012 \cite{hennessy2019}.
The sentence that usually follows, ``so the industry had to go multicore'', does not
follow. Physics forbade a single processor getting faster at constant power and said
nothing about the response; accelerators, better software or accepting the performance on
offer were all available. Multicore was a choice, and reporting a chosen response as a
physical necessity is the error this paper exists to catch. The four bills are the floor
of the space in which engineering happens; the rest of the space is what the model left
out.

% ===========================================================================
\section{Eight Idealisations}\label{sec:idealisations}
% ===========================================================================
Take the model's assumptions one by one and ask what a real machine pays that the model
does not. Eight answers come back, and those are the \emph{idealisations}; to pay under
an idealisation is to pay the cost it hides. They sort into three groups. The first five
are physical. The sixth sits at the boundary between the program and the people who write
it. The last two come from the world outside it. The two censuses that follow the list
show where formal knowledge stops.

\textbf{1. The tape is unbounded and free.} A real machine is finite and a byte has a
price, which has fallen about six orders of magnitude since 1987 and is the input to
Gray's formula in \S\ref{sec:practice}. Generates the storage tiers as a price ladder,
what is kept and what is recomputed, and the rule to keep the working set in memory.

\textbf{2. Movement is free.} False by the latency bound, and measured at every distance:
sequential memory about half a nanosecond, random memory twenty, local flash a hundred
microseconds, disk ten milliseconds, object storage eighty \cite{napkinmath}, eight orders
of magnitude, so that if the fastest were one second the slowest would be three years;
and within the processor a last-level cache miss costs about 250 core cycles on Skylake
and 420 on Zen~4 \cite{sevencpu,chipsandcheese}, the product of memory latency improving at
about 7\% a year against processor speed at 25 to 50\% \cite{hennessy2019}, foreseen in
1994 under the subtitle ``implications of the obvious'' \cite{wulf1995}. Generates the
memory hierarchy, caching, indexing, locality as a first-class concern, columnar storage,
data-oriented design, and the rules every engineer knows: index the filtered column, cache
what is expensive and stable.

\textbf{3. There is one head.} Real machines have many, because a machine near its rate
limit must be parallel and machines far below it went parallel by choice
(\S\ref{sec:floors}), and the heads must be coordinated. Measured, indirect costs dominate: a
context switch 3.8 microseconds directly and hundreds through cache pollution
\cite{li2007}; a system call 150 cycles directly and about 12,300 in full footprint
\cite{soares2010}. Generates concurrency control, transactions, isolation levels, memory
models.

\textbf{4. One tape, in one place, that does not fail.} Failure is the third bill.
Distance: an in-process call about a nanosecond, a system call three hundred, a
same-region network call 250 microseconds \cite{napkinmath}, or two to four with kernel
bypass \cite{kalia2019}; three to eight orders of magnitude depending on the
infrastructure.
Generates distribution, protocols,
replication, consensus.

\textbf{5. The machine is alone.} Idealisation 3 is many heads the designer controls; this
is heads the designer does not: other tenants, other authors, attackers. An attacker is
somebody who reads a variable the designer declared irrelevant: Spectre read discarded
timing, Meltdown read state the architecture declared invisible. So what changes here over
time is not a price. It is whether the premise that some variable is irrelevant is still
true. Generates isolation, capabilities, sandboxing, tenancy.

\textbf{6. The program is given.} The sheet is written by people of limited memory and
attention and revised by others for as long as the system matters, and nobody holds the
whole of it. Simon's argument that anything complex assembled in interruptible steps ends
up a nearly decomposable hierarchy is the founding statement \cite{simon1962}. Generates
languages as tools for people, abstraction, modularity, interfaces, version control,
review, testing, architecture; nearly every rule in a style guide is a price paid for this
cost.

\textbf{7. The input is given.} What a system should do lives in people's heads,
incompletely, and is discovered while building; this is Naur's programming as theory
building \cite{naur1985}. Generates requirements, product work, every method that
shortens the loop between building and finding out whether it was wanted.

\textbf{8. The job is fixed and finishes.} Maintained systems never halt; they are changed
while running, the world moves under them, and their structure degrades unless effort is
spent against it \cite{lehman1980}. Generates maintenance, operations, observability,
migration, the economics of technical debt.

The machine in this list is not any particular modern computer. A processor, a phone, a
data centre and a fleet are each Turing's machine, matter arranged so that its state
changes correspond to the five operations for every input, and each has a different price
list under the same eight; a new substrate changes the prices, never the list, and the
walls of \S\ref{sec:cannot} do not move for any of them. Computation that never passes
through symbols, the analogue kind, is set aside here as Turing set aside the choice
machine.

The bills set floors under the physical five without pricing each one: the latency bound
under 2, stepping under 3 and, as the speed of light, under 4, failing under 4; erasing
binds nothing. The first has a market price and no floor that binds. The fifth has
measured overheads, and behind it stands a premise rather than a price.

\subsection{Two Censuses}

Two counts, with different instruments, find where formal knowledge runs out. The first
count is of cost models. The physical five have them, four with published constants and
the fifth with measured overheads for isolation; the other three have none. The second
count is of theorems, each sorted under the idealisation whose cost its question
presupposes, so that a result about checking a sheet somebody wrote goes under 6 and a
result about a tape in more than one place goes under 4. The durable results of computing
pile up under 6 (undecidability, Rice, Myhill and Nerode, Church and Rosser, Codd, Hindley
and Milner, Cook and Levin) and under 4 (FLP, CAP, Paxos, Raft, replicated data types).
Under 5 the unconditional results are about what a foreign head can learn from what it
reads, not about keeping it off the tape, and they are few. Under 7 the results are about
learning functions from samples, not about eliciting intent from people. Under 8 there are
none, since Lehman's laws are regularities, Conway's is an observation, and Brooks's is an
argument.

The two world idealisations were enumerated, as interaction with an outside world, a
program that changes and a machine that never stops, by van Leeuwen and Wiedermann
\cite{vanleeuwen2001}. This paper claims the full list, the sorting, and that the sorting
predicts where formal knowledge runs out. Completeness is a failed search, refutable as
Turing's thesis is: exhibit a ninth cost the model discards that is not a combination of
these.

The list fixes the line between the interpreter's first two tiers. Idealisations 7 and 8
are the world, what people want and what they will want next; everything from 6 down is
the machine and the people who write for it. The first tier works beyond that line, no
cost model reaches it, and that is why the paper leaves it out. The idealisations say what
a real machine pays. They do not yet say where, and where is the next question.

% ===========================================================================
\section{Rows}\label{sec:rungs}
% ===========================================================================
A \emph{row} is the level of a system at which an operation is performed and paid for,
from the machine to the institution. The five operations recur at every row, at prices
separated by many orders of magnitude, and the rest of the paper refers to rows by name.

\begin{center}
\small
\begin{longtable}{@{}p{0.13\textwidth}p{0.14\textwidth}p{0.14\textwidth}p{0.14\textwidth}p{0.13\textwidth}p{0.16\textwidth}@{}}
\toprule
\textbf{Row} & \textbf{Read} & \textbf{Write} & \textbf{Move} & \textbf{State} &
\textbf{Halt} \\
\midrule
\endhead
Machine & load & store & address & register & halt instruction \\
\addlinespace
Language & evaluate & assign & control flow & scope & return \\
\addlinespace
Program & call a getter & call a setter & follow a reference & object state & terminate \\
\addlinespace
Process & receive & send & schedule & process state & exit \\
\addlinespace
System & query & commit & route & session & shut down \\
\addlinespace
Organisation & ask someone & tell someone & hand off & ownership & disband the team \\
\addlinespace
Institution & read the spec & publish a version & migrate & the standard & deprecate \\
\bottomrule
\end{longtable}
\end{center}

Read the \emph{read} column down with prices: a register in a fraction of a nanosecond,
main memory in tens of nanoseconds, flash in a hundred microseconds, another continent in
tens of milliseconds. Then past the machine, where a read means asking a colleague: minutes
at the next desk, days in another team, unavailable if they have left. That last step makes
Conway's law a cost law: the organisation's communication graph is the price list for
reads at that row \cite{conway1968}.

The \emph{halt} column matters more. Stopping something is one instruction at the machine
row, a restart at the process row, a deployment with a rollback plan at the system row, a
reorganisation at the organisation row, a deprecation announced years ahead at the
institution row. \emph{The cost of stopping rises with the row at which the decision was bound.} That one
observation is the general form of much architectural advice: do not bake configuration
into code, a business rule into a migration, a policy into a protocol, a team structure
into a public API. Each of these says that a decision was bound at a row whose halt cost
exceeds what it was worth, and each becomes checkable once the row is named. Many of the
field's loudest arguments are two rows mistaken for one.

\subsection{Where a Quantity Is Preserved}\label{sec:whopays}

Because every operation exists at several rows, anything that must be preserved could be
preserved at any of them, and choosing the row is a decision. The end-to-end argument
addressed exactly this, and Saltzer, Reed and Clark stated the under-determination: the argument
``does not tell us where to put the early checks'', it is ``a guideline'', and they supplied
the cost model that decides, with a crossover in file size beyond which end-to-end retry is
wrong \cite{saltzer1984}. That is a rule published with its conditions.

\begin{center}
\small
\begin{tabular}{@{}p{0.28\textwidth}p{0.66\textwidth}@{}}
\toprule
\textbf{Preserved quantity} & \textbf{Rows that could carry it} \\
\midrule
Uniqueness of a key & the database constraint, the data-access layer, the service, the client \\
Retry on failure & the transport, the client library, the application handler, the human \\
Type correctness & the compiler, the runtime, the schema registry, the test suite \\
Authorisation & the gateway, each service, the database, the query \\
Idempotency & the caller, the broker, the handler, the store \\
Ordering & the log, the queue, the consumer, the reconciliation job \\
\bottomrule
\end{tabular}
\end{center}

\emph{An architectural decision is an assignment of a preserved quantity to a row, and its
cost is the halt cost of that row.} That is one kind of decision. Whether every decision has this shape is the next question.

% ===========================================================================
\section{Decisions}\label{sec:decisions}
% ===========================================================================
A decision exists only where two alternatives differ in cost. That is why nothing is
ever chosen in Turing's machine. With every cost set to zero it does not matter where a
write goes or how far away a cell is; nothing races, nothing is lost, nobody else is
present, nobody has to be asked, nothing has to be maintained. So the paper need not list decisions, which would be
hopeless; it inherits the enumeration from the idealisations.

\begin{quote}
\emph{Every decision in software engineering is a choice between alternatives that differ in
what they pay under one or more of the eight idealisations, and a decision is complicated in
proportion to how many it pays under at once.}
\end{quote}



\subsection{A Grammar of Decisions}

Three coordinates suffice, and all three are already in the paper. The \emph{operation} is
a read, a write, a move, a state change or a halt, since those are the only things a
machine does. The \emph{row} sets the halt cost, which is what a decision costs when it is
wrong. The \emph{idealisation set} is the set of idealisations under which the alternatives
differ; it is what the decision is answerable to, and it is the coordinate nobody records.
A decision is a point in operation $\times$ row $\times$ idealisation set, and its
\emph{shape} is which point it occupies.

\begin{center}
\small
\begin{longtable}{@{}p{0.30\textwidth}p{0.14\textwidth}p{0.20\textwidth}p{0.24\textwidth}@{}}
\toprule
\textbf{Decision as people say it} & \textbf{Operation} & \textbf{Row} & \textbf{Paid
under} \\
\midrule
\endhead
Index the column the query filters on & read & system & 1, 2 \\
Keep the working set in memory & read & system & 1, 2 \\
Make it a microservice & move & system, organisation & 2, 3, 4, 6, 8 \\
Add a boolean flag & state & program & 6, 8 \\
Pin the dependency version & write (by another head) & program, institution & 5, 6, 8 \\
Write a test for it & halt (a gate) & program & 4, 6, 8 \\
Ask the founder before building & read & organisation & 7 \\
Put it behind a feature flag & halt & system & 8 \\
Retry on failure & move & any of four rows & 4 \\
Rename this for clarity & write & language & 6 \\
\bottomrule
\end{longtable}
\end{center}

The decisions everyone calls hard pay under many idealisations at once, and microservices,
the worst offender, is five debates wearing one name, which is why it never ends. The
decisions everyone calls trivial are often those whose idealisation set was misidentified,
which is where the deceptively simple failures come from. The rest of the felt
complexity is a row nobody named or an idealisation set that moved with the prices, and
naming the coordinates recovers both.

\subsection{A Dependency Failure}\label{sec:dependency}

A library ships a new version and the build breaks, or passes and production breaks. It
looks like a lookup, resolved by reading a changelog, and it is not. A dependency is a
region of the program's tape that a head its author does not control may write to, and
the program's correctness rests on those symbols staying what they were. That is
idealisation 5. It is composed with 6, since the sheet was written by people who did not
know about this program, and with 8, since the world moves while the machine runs. ``Keep dependencies up to date'' addresses 6 alone, and the failure comes from the other two.
Once that is visible, the resolution is three questions. Which symbols in the foreign region does the program actually read? That is the
\emph{cut}; semantic versioning is an attempt to write it down as a promise, and semver
failures are the promise being narrower than the cut. At which row was the version bound?
In a lockfile the halt cost is a one-line revert; in every service's base image, a
coordinated redeploy; in a platform the team does not control, a deprecation cycle it does
not own. Who else has a head over the region? Two dependencies reading incompatible symbols
from a third have no fix at the program row; that is why diamond dependencies feel
different. It is the same problem, one row up. \emph{A dependency failure is a write by a foreign
head to a shared region whose cut was never written down or written down wrongly, and its
cost is the halt cost of the row at which the version was bound.} Pin, float, vendor,
lockfile, monorepo and reproducible builds are choices of that row.

\subsection{The Census}\label{sec:census}

Decisions leave traces in
pull requests, incident reports, design reviews and migration write-ups. Take a sample
from a codebase's history, code each by operation, row and idealisation set, and record
whether coding was possible. Decisions that code onto one idealisation are the simple ones. Those that code onto several are the composed ones. Those that do not code at all
are either a ninth idealisation or a coding error, and either way they are the interesting
residue. It is days of work for any team whose history is written down, and it is Prediction~1.
A decision with a shape can be judged, and judging it means judging the rule it was made
with.

% ===========================================================================
\section{The Test}\label{sec:practice}
% ===========================================================================
The list of things a rule can be answerable to is now complete. A rule can answer to a wall, to a bill, to a cost under an idealisation at some
row, to a measurement, or to nothing but the fact that others follow it. That list makes
one question askable of any imperative.

\begin{test}[What is this practice answerable to?]\label{test}
For any practice stated as an imperative, name which of the following it answers to, and
the specific instance:
\begin{enumerate}[leftmargin=1.6em,itemsep=1pt]
\item[\textbf{(a)}] \textbf{A cost.} Under which of the eight idealisations, paid at which
row, with what measured ratio and threshold; where physics sends a bill, which bill.
\item[\textbf{(b)}] \textbf{A theorem.} Which one, with its fine print, and whether it
licenses the practice or merely constrains it.
\item[\textbf{(c)}] \textbf{A measured regularity.} Which study, on what population, when,
and whether replicated.
\item[\textbf{(d)}] \textbf{Uptake.} Nothing outside the population that adopted it; its
value is coordination and its truth condition is that others follow it.
\end{enumerate}
A practice offered with no answer is an assertion that something is optimal with no
statement of what it is optimal for, and there is nothing in it to check.
\end{test}

A verdict of (d) is not a condemnation: driving on the left answers to nothing but the fact
that others do, and it is among the most valuable rules there are; what (d) forbids is
defending a convention as a discovery. Vincenti established that a
whole category of engineering knowledge is prescription without theoretical grounding
\cite{vincenti1990,abran2006}. What he could not say, having no physical layer to check
against, is which prescriptions have a cost underneath them. Sections \ref{sec:cannot} to \ref{sec:idealisations} are that physical layer.

\subsection{Four Rules}\label{sec:verdicts}

Four practices, chosen because their answers differ.

\textbf{Keep the working set in memory.} Verdict (a): the latency bound, paid at the system
row, and quantified because Gray and Putzolu wrote the conditions down. In 1987
Gray and Putzolu derived the break-even interval for keeping a page in memory as
$A\$/(M\$\cdot B)$, the price of a disk access per second over the price of a byte of
memory times the page size. Their inputs gave 400 seconds, hence the five-minute rule
\cite{gray1987}. In the same paper is the sentence this paper is about:

\qt{``These results depend on current price ratios of processors, memory and disc accesses.
These ratios are changing and hence the constants in the rules are changing.''
\cite{gray1987}}

Because the formula was public, everybody could watch it expire. In 1997 it gave 266
seconds, still five minutes only because page size had grown eightfold and absorbed a price
change \cite{graygraefe1997}. In 2007 flash split the question into three answers
\cite{graefe2007}. In 2019 memory against disk reached four hours and flash against disk a
day, mechanism stated \cite{appuswamy2019}. A 2025 preprint that charges for processor
time reports seconds \cite{zhang2025}.  This is a fully anchored practice, and there are very few.

\textbf{Do not optimise prematurely.}\label{sec:knuth} Verdict (a), with the regime lapsed.
Since the subject is dropped conditions, here is the whole paragraph:

\qt{``There is no doubt that the grail of efficiency leads to abuse. Programmers waste
enormous amounts of time thinking about, or worrying about, the speed of noncritical parts
of their programs, and these attempts at efficiency actually have a strong negative impact
when debugging and maintenance are considered. We \emph{should} forget about small
efficiencies, say about 97\% of the time: premature optimization is the root of all evil. Yet
we should not pass up our opportunities in that critical 3\%. A good programmer will not be
lulled into complacency by such reasoning, he will be wise to look carefully at the critical
code; but only after that code has been identified.'' \cite{knuth1974}}

Only the answer travelled. The cost it answered, debugging and maintenance time, is in the
first two sentences and did not travel, which is why the proverb now supports positions
Knuth argued against. The regime lapsed too. The procedure, identify the critical code by
measuring and only then optimise it, assumes running cost is dominated by instruction
count. Since then a cache miss has become 250 to 420 cycles
\cite{sevencpu,chipsandcheese} and a memory fetch ten thousand times the energy of an add
\cite{dally2022}, so the expensive decision is now made when the data structure is chosen,
before there is anything to profile \cite{acton2014}. The principle survives; the procedure
does not.

\textbf{Keep lines to eighty columns.} Verdict (d). Hollerith's 1890 card was sized to fit
boxes made for the Treasury Department \cite{jones-cards}. IBM's 1928 rectangular holes
packed eighty columns into that width \cite{ibm-card}. Terminals inherited the number, and
mid-1970s terminals shipped at many other widths, so no display forced it
\cite{shirriff2019}. The numbers have since drifted, PEP~8 at 79, Google's Java guide at
100, Black at 88, Go at none, and the practice is still enforced. Reading research on
prose finds that longer lines read faster and that preference disagrees with performance
\cite{dyson2004,bernard2002,nanavati2005}. On code there is nothing; the nearest work trains
on readability ratings, which is circular \cite{buse2010}. The rule is a convention, and a
line limit is worth having; what the test exposes is that nobody has measured which limit,
and that the one in use is the width of a card.

\textbf{Limit work in progress.} Verdict (b), and the theorem does not license the
practice. Little's law, $L=\lambda W$, where $L$ is the work in progress, $\lambda$ the throughput
and $W$ the time each item spends in the system, holds under almost no hypotheses
\cite{little1961,stidham1974}, so the relationship is not negotiable. But it is an
identity, not a mechanism. Cut $L$, and the identity is equally consistent with $W$
falling at constant $\lambda$, which the advocacy promises, and with $\lambda$ falling at
constant $W$, which nobody advertises. The practice's whole empirical content sits in a
constant-throughput clause the theorem says nothing about, and that is why the practice
literature's own assumption lists are stronger than the theorem's
\cite{anderson2016,vacanti2015}. The mechanism that would justify it, Kingman's
heavy-traffic result \cite{kingman1961}, carries exactly the assumptions software
development violates. The one empirical study, of 8,000 work items, called for replication
and has not had it \cite{sjoberg2018}.

\begin{center}
\small
\begin{longtable}{@{}p{0.26\textwidth}p{0.15\textwidth}p{0.51\textwidth}@{}}
\toprule
\textbf{Practice} & \textbf{Verdict} & \textbf{What it is answerable to} \\
\midrule
\endhead
Keep the working set in memory & (a) & The latency bound; a formula published in 1987 and
recomputed four times since \\
\addlinespace
Do not optimise prematurely & (a), lapsed & Instruction count against maintenance time; the
denominator changed \\
\addlinespace
Eighty-column limit & (d) & A card whose width a Treasury box fixed in 1890; nothing measured
on code \\
\addlinespace
Limit work in progress & (b), insufficient & Little's law, an identity that cannot license
the intervention \\
\bottomrule
\end{longtable}
\end{center}



\subsection{Contradictory Proverbs}\label{sec:pairs}

The field's proverbs come in pairs that seem to contradict, and do not. Each pair is one
inequality seen from both sides: a case analysis whose separating condition was dropped
from both branches.

\begin{center}
\small
\begin{longtable}{@{}p{0.24\textwidth}p{0.34\textwidth}p{0.36\textwidth}@{}}
\toprule
\textbf{Pair} & \textbf{The two costs} & \textbf{Which side a system is on} \\
\midrule
\endhead
DRY against decouple & A change made twice, against a boundary coupling two sources of
change & Do the two copies change together or separately in the history \\
\addlinespace
YAGNI against design for change & Unused flexibility, against retrofitting it & How likely
is the change, and what does reversal cost \\
\addlinespace
Small services against modular monolith & Coordination cost of a shared deployable, against
operational and latency cost of a boundary & Releases last quarter that needed two teams;
fraction of tail latency in remote calls \\
\addlinespace
Normalise against denormalise & Write and consistency cost, against read and join cost &
Read-to-write ratio, and whether reads span the join \\
\addlinespace
Fail fast against be liberal & A wrong result, against a refused request & Who bears the
failure, and is the input under the team's control \\
\addlinespace
Explicit against convention & Stating it, against a reader inferring it & How many readers,
how expert, how often \\
\addlinespace
Pin against float & Reversal cost of a foreign write, against staleness and security debt &
How often the region read from changed across the cut \\
\bottomrule
\end{longtable}
\end{center}

The third column is instructions, and most of them can be measured on a given system
within a week. Put the costs back and the principles of software engineering come out,
each one the cheaper side of an inequality between two costs, so a unit of engineering
knowledge has a normal form: \emph{for decisions of shape (operation, row, idealisation
set), alternative A is cheaper than B when this inequality between these two measurable
quantities holds, and the observable that tells which side a system is on is this}. A
theorem is the case where the inequality holds in every regime. A lookup is the case where
there is no alternative. A convention is the case where the only cost on either side is
the cost of differing from everyone else. Two entries of the third column cannot be
measured within a week: who bears a failure is a judgment about people, and the
likelihood of a future change is a different kind of quantity, the subject of
\S\ref{sec:stops}. Which leaves the question of why the conditions were ever
dropped.

% ===========================================================================
\section{How Conditions Are Lost}\label{sec:whykeydropped}
% ===========================================================================
Conditions are lost by four mechanisms. None is a failure of character, and they compound.
The first is skill acquisition, which compiles declarative knowledge into procedures that
fire on a pattern, so that the reason is no longer retrievable from the procedure
\cite{anderson1982}. The second is that expert decision is
recognition, not comparison. Fireground commanders compared options in fewer than twelve
percent of 156 decision points \cite{klein1986}, and software designers consider fewer
options the more structured the decision \cite{zannier2007}, so the cost model that once
made the action cheapest never enters working memory. The third is that making context
explicit is expensive, and the person who pays is not the person who benefits
\cite{shipman1999,grudin1996}. The fourth is formalisation itself: turning an argument into an equation keeps the
variables and drops the constants and the setting. Amdahl's 1967 paper has
no equations and one empirical constant, data management overhead near forty percent of
instructions \cite{amdahl1967}. The algebra everyone calls Amdahl's law kept the equation,
dropped the constant and silently fixed the problem size, and Gustafson needed a page and
a half, twenty-one years later, to put the condition back \cite{gustafson1988}. The latency bound of \S\ref{sec:bills}
has travelled the same way: the classical model justifies its cost function by circuit
delay \cite{aggarwal1987}, the textbook by locality and manufacturing \cite{hpqa}, and the
cube-root claim circulates in forum posts that cite nothing. Together these predict practitioners holding beliefs not indexed to their own projects' conditions,
and a survey of 564 Microsoft engineers found exactly that \cite{devanbu2016}.

The general form is old and belongs to economics: QWERTY is a practice optimal for a lapsed regime,
retained by convention and network effects \cite{david1985,arthur1989}, and inside software
the same thing is the vaporisation of architectural knowledge \cite{jansenbosch2005}. This
paper adds only the expression of the cached rule as a minimisation over a stated
cost model, so that the regime a rule assumed becomes a quantity and a threshold rather than a story.

If that is right, a practice should invert when its inequality flips, with
the same kind of people giving the same kind of reason in the opposite direction.
Migrations are written up in public, so this can be checked.

\subsection{Inversions}\label{sec:inversions}

\begin{center}
\small
\begin{longtable}{@{}p{0.20\textwidth}p{0.36\textwidth}p{0.36\textwidth}@{}}
\toprule
\textbf{Case} & \textbf{The ratio that moved} & \textbf{What the practitioners said} \\
\midrule
\endhead
Store against recompute & Instruction cost against memory stall & Acton 2014, with the
cache-line accounting \cite{acton2014} \\
\addlinespace
Buy hardware against optimise & Free performance against engineering cost; Dennard ended
2007 to 2012 & Python to hand-tuned vector code is 62,806$\times$ on one kernel
\cite{leiserson2020} \\
\addlinespace
Normalise, denormalise, back & Developer time against machine throughput, twice in opposite
directions & Codd's opening obligation \cite{codd1970}; Stonebraker's 900 against 160,000
messages per second \cite{stonebraker2005}; F1 calling eventual consistency ``an
unacceptable burden to place on developers'' \cite{shute2013} \\
\addlinespace
Monolith, services, monolith & Coordination cost of a shared deployable against operational
and latency cost of a boundary & Segment \cite{noonan2018}; Istio \cite{mendonca2021}; Uber,
fifty services across twelve teams for one root cause \cite{gluck2020} \\
\addlinespace
Threads against processes & Argued to have crossed; transition cost priced higher & Lane:
``I think this will be a disaster\dots largely silently'' \cite{lane2023} \\
\addlinespace
Site isolation & Nothing moved; a premise was falsified & After Spectre ``it is no longer
safe to render documents from different web sites in the same process'' \cite{reis2019} \\
\addlinespace
Cloud and back & Rent against own, conditions stated by the practitioner & Heinemeier
Hansson: \$3.2m budget, \$1.5m projected saving, and the regime under which renting stays
right \cite{dhh2023} \\
\bottomrule
\end{longtable}
\end{center}

Three things the table cannot carry. Site isolation bounds the framework, because inversions of that
kind come from outside any cost model. And two counter-cases. Prechelt and colleagues found a design pattern's flexibility probably worth
having absent a clear reason otherwise, and could identify the exceptions only by
``software engineering common sense'' \cite{prechelt2001}; some cached answers are robust
across a regime no system leaves. Menzies and colleagues found no evidence across 171
projects that defects get more expensive to fix the longer they linger \cite{menzies2017},
a belief that has justified a great deal of process. The table records inequalities that
have already flipped; the ratios still moving say which flip next.

\subsection{Hypotheses}\label{sec:forward}

Each hypothesis below names the ratio that moves, what should invert, what should not, and
an observable.

\begin{hyp}[Cheap memory splits the memory practices in two]
Memory cost per byte has fallen about six orders of magnitude since 1987, so every
practice anchored on the price of a byte inverts. But the bound is on distance and more
memory is a larger $x$, so practices anchored on the distance to a byte, locality,
sequential access, data-oriented layout, intensify. \emph{Observable}: the ratio of slowest
to fastest in-memory access on a machine across generations; flat or falling while
capacity rises refutes this. Half the evidence is in: the 2019 recomputation pushed the
byte-price crossover to four hours \cite{appuswamy2019}; the distance half awaits the
observable.
\end{hyp}

\begin{hyp}[Returns concentrate in packaging and near-data computation]
$\Omega(x^{1/d})$ cannot be bought off, but it has two letters that can be moved. Raise
$d$ by stacking, which changes the exponent from a half to a third; at a billion cells
that is worth a factor near thirty, and there is no fourth dimension after it. Or lower $x$
by moving computation to the data.
\emph{Observable}: the share of a hardware generation's improvement attributable to
organisation and placement rather than device speed.
\end{hyp}

\begin{hyp}[Distrust practices track fleet size rather than device quality]
Per-device error rates fall every generation and the practices built on distrust have not
relaxed in thirty years, because the governing quantity is rate times count and the count
has risen at least as fast. They relax first where fleets stop growing: single machines,
edge and embedded work. \emph{Observable}: rate times count for the system actually run,
now and five years ago.
\end{hyp}

\begin{hyp}[Practices re-anchor from time to energy]
The Landauer floor limits nobody; heat removal limits everybody and stopped improving when
Dennard scaling ended. Rules justified by latency will be restated with energy as the
objective, and the restatement changes answers about idling, batching and speculation.
\emph{Observable}: joules per unit of work on the hot path, which most teams do not
measure.
\end{hyp}

\begin{hyp}[Gap practices expire and floor practices do not]
Between the fibre floor and delivered pings there is a factor of about three to win.
Practices that hide the gap, prefetching, connection reuse, edge caching, work against a
finite shrinking margin; practices that respect the floor, partitioning by geography,
asynchronous defaults, a deliberate consistency model, are permanent. \emph{Observable}: a
system's median round trip against the great-circle fibre floor for the same endpoints.
\end{hyp}

\begin{hyp}[The binding constraint moves from production to reversal]
The cost behind idealisation 6 has no floor, which is why it has fallen as far and as fast
as it has. A large part of good practice answers to that cost and to nothing else: do not
repeat yourself, abstraction to avoid retyping, style guides, small reviewable changes,
``we cannot afford to rewrite it''. The halt column does not fall with it. So as writing
gets cheap, the practices that survive are the ones that minimise the row at which
decisions bind, not the amount of code. DRY weakens. YAGNI strengthens. Feature flags and
late binding rise in value against anything whose benefit was saved typing. \emph{Observable}: the fraction of a team's change
cost attributable to writing, against reversal and verification, across its last twenty
changes. If writing collapses and practice does not reorganise around reversal, either
those practices never answered to production cost or the mechanism is wrong.
\end{hyp}

The last hypothesis is the hinge, because the cost it names is the one now being removed.

% ===========================================================================
\section{The Cost of a Decision}\label{sec:formula}
% ===========================================================================
What happens when a cost is removed can be computed, because a decision has a cost with
four terms. For a chosen alternative $a$,
\[
C(a) \;=\; P(a) \;+\; V(a) \;+\; p_{\text{wrong}}(a)\cdot H(\text{row}(a)) \;+\; I(a).
\]
Here $P$ is production, what it takes to write the change. $V$ is verification, what it takes
to establish that the change does what was intended. $p_{\text{wrong}}$ is the chance the
decision is wrong given what is known when it is made, and $H$ is the halt cost of the row
it binds to, so their product is the expected cost of undoing it. $I$ is information, what
is spent reducing $p_{\text{wrong}}$ before committing. The interpreter picks the $a$ that
minimises $C$, and the field's history reads as a record of which term dominated. For most
of that history $P$ was large. Where $p_{\text{wrong}}$ and $H$ were both high, practices
minimised their product: feature flags, canaries and staged rollouts bind a decision
temporarily at a row with lower $H$ while $p_{\text{wrong}}$ is still being reduced. The formula also says when to stop and
ask somebody: when the reduction in $p_{\text{wrong}}\cdot H$ that the answer would bring
exceeds $I$, the cost of getting it; and $I$ depends on who is asked, so it also says whom,
the cheapest holder of the answer.  

\subsection{Production Goes to Zero}\label{sec:pcollapse}

Set $P$ to nearly zero, as a capable code generator does. Every practice
that was minimising $P$ stops being answerable to anything. Any decision whose
idealisation set was $\{6\}$ alone loses its trade-off, and duplication, boilerplate and
much of what style guides exist for fall in that class. What survives pays under something
other than 6. Idealisations 1 to 4 are physics at the machine rows and do not move. Number
5 gets worse, because more heads are writing. Numbers 7 and 8 are untouched, because a
model that writes cheaply does not know what the founder meant or what the world will want
next year. What is left is
\[
C(a) \;\approx\; V(a) \;+\; p_{\text{wrong}}(a)\cdot H(\text{row}(a)) \;+\; I(a),
\]
every term about being wrong rather than producing. This is the second tier of
\S\ref{sec:tiers} becoming visible: the software intention was always what was being
decided, and the code was the receipt.

\subsection{Verifying Decisions}\label{sec:verification}

Current verification assumes $P$ is large: a person writes slowly and a reviewer reads what
was written. It checks the artefact. That made sense when the artefact was the expensive thing, and it
stops making sense when the artefact is cheap and regenerable, because reading is now the
expensive step and it verifies the wrong object. What needs verifying is the decision, and
\S\ref{sec:decisions} says what one is, so there are four checks. The \emph{idealisation set} was identified correctly, so that the trade-off
priced is the one the world charges for; the dependency failure of \S\ref{sec:dependency}
fails this check. The \emph{inequality's inputs were measured} on this system rather than
carried over from someone else's. The \emph{row is the lowest} that satisfies the
constraint, since binding higher buys reversal cost that was not needed. And the
\emph{reversal path has been exercised}, because a rollback that has never run is a hope.

The checks split along the idealisations. Measured inputs and a working reversal path are
checked by machines against machines. Whether the idealisation set is right is mechanical
for 1 to 4, which have cost models. For 5 and 6 it is one head checking another's sheet
and another's premises, which a machine does without tiring at the thousandth call site.
For 7 and 8, which concern what the world wants, it is not mechanical at all. Whether the
row is right depends on who depends on the binding. So the physical idealisations are
verified by machines against measurements, the boundary idealisation by machines against
each other, and the world idealisations by whoever owns the goal, who need not be an
engineer: finance, for whether invoices may be altered; compliance, for whether data may
leave the region; the founder, for simplicity now against flexibility later.  

Each moment a machine doing engineering work stops to ask a person can be sorted under
an idealisation in the same way. Information that exists but is expensive to reach
at the organisation row is 2. An objective nobody has decided is 7, the residue that can
only be routed, never automated. An unknown dependency is 5. A missing abstraction is 6,
an ordinary engineering problem. Over thousands of interventions the residue, the stops that remain, is then not a
mystery called judgment but a distribution over eight categories with a known direction:
1 to 4 already machine-checkable, 5 and 6 shrinking with tooling, 7 and 8 where the engineer stays. That is Prediction~7.

\subsection{The Record}\label{sec:record}

Both the verification above and the test need a written record of a rule with its
conditions attached, and no existing format holds one. The pattern form asks for context
and forces, qualitatively \cite{alexander1977,gamma1994}. Architecture decision records ask,
in one line, that changes of context be noted \cite{nygard2011}. Neither asks for a ratio,
a threshold or a date.

\begin{center}
\small
\begin{longtable}{@{}p{0.20\textwidth}p{0.74\textwidth}@{}}
\toprule
\textbf{Field} & \textbf{What goes in it} \\
\midrule
\endhead
\textbf{Rule} & The imperative, as people say it. \\
\addlinespace
\textbf{Shape} & Operation, row, idealisation set. \\
\addlinespace
\textbf{Answerable to} & (a) cost, (b) theorem, (c) measured regularity, or (d) uptake,
with the instance named. \\
\addlinespace
\textbf{The two costs} & Named as measurable quantities, not virtues. \\
\addlinespace
\textbf{Inequality} & The condition under which the rule is the cheaper option. \\
\addlinespace
\textbf{Crossover} & The value at which it stops being cheaper, with the inputs used. \\
\addlinespace
\textbf{Observable} & What is measured on \emph{this} system to know which side it is on. \\
\addlinespace
\textbf{Reversal} & What undoing costs at the bound row, and whether the path has been
exercised. \\
\addlinespace
\textbf{Inversion} & What the rule becomes on the other side. Inability to state it means a
theorem, a regime wider than anything the holder will build, or an unexamined belief. \\
\addlinespace
\textbf{Last recomputed} & A date, because a verdict can change when a price ratio
crosses, and a rule with no date is a rule nobody has checked.
\\
\bottomrule
\end{longtable}
\end{center}

The record serves two readers. A person following a rule
cannot usually say when it would be wrong. The record writes the regime beside the rule,
with a date, so that following and checking become separate activities. A machine told to
prefer small services can only pattern-match. Given the record, it can measure the named
observable and say which side of the inequality the system is on, and that matters most
where the advice is wrong. And the machine hands back the record, and a person verifies
it when the code is no longer worth reading. That is as far as the record goes, and the limit is exact.

% ===========================================================================
\section{The Limit}\label{sec:stops}
% ===========================================================================
A practice resting on a physical cost can be recorded and recomputed, because the cost has
a model. The other three idealisations have no cost models, 7 and 8 have no theorems, and
nobody has written down what has to be preserved there. That is a fact about the type of
the quantity, not a failure of effort. Every episode in which generation became mechanical
has had the same shape, and it is the shape \S\ref{sec:knowledge} predicted: somebody made
a piece of reality hold still by writing down what had to be preserved, and generation
followed.  The C standard's ``(only) the observable behavior of the abstract machine'' made
optimisation a discipline \cite{c11}. Serial equivalence made concurrency control
mechanical \cite{eswaran1976}. Relational algebra made query planning mechanical, and Codd
opened by stating the preserved quantity as an obligation to users \cite{codd1970}. The
five-minute rule is the same fact where reality does not hold still: the quantity was
written down, the prices moved, and the answer was recomputed each time.

Architecture is where this never happened, because its preserved quantity is a bet about
the future. Parnas said so in 1972: begin with a list of the design decisions \emph{likely
to change}, and hide each in a module \cite{parnas1972}. The list is not derived from the
code. It is a person's guess about what the world will ask for next. Two companies holding
byte-identical code, one serving a process unchanged for a decade and the other a market
that reinvents itself quarterly, need different decompositions of the same bytes, and
nothing computed from the bytes says which decomposition is right.

\begin{center}
\small
\begin{tabular}{@{}p{0.22\textwidth}p{0.30\textwidth}p{0.40\textwidth}@{}}
\toprule
\textbf{Object} & \textbf{A formula over} & \textbf{Where the inputs live} \\
\midrule
A program & the tape & in the text; fully present \\
An optimisation & the machine model & in a published cost model \\
A practice & two costs and a ratio & in nobody's document, but derivable by measurement \\
An architecture & a distribution over future requests & in nobody's document, and not
derivable by measurement \\
\bottomrule
\end{tabular}
\end{center}

The last row is why the engineer stays at idealisations 7 and 8.

% ===========================================================================
\section{Rival Explanations}\label{sec:alternatives}
% ===========================================================================
The account is now complete, and it is worth holding only if the rival accounts fail.
Each of the four below is partly right, and this paper keeps the part that is.
What none of them does is say where a rule stops being right.

Suppose first that the rules of the field are conventions all the way down: index the
filtered column because others do, and for no other reason. Then a rule would change only
when the people changed, and it would track no quantity. But the practices of
\S\ref{sec:inversions} reversed in public, several of them in both directions, with the
practitioners naming the number that had moved. And where a formula was published, the
number was recomputed four times as the prices moved while the rule's name stayed the
same. Conventions exist, and the test has a category for them. A convention, though, has
no number to track, and these tracked one.

Suppose the opposite: that engineering is applied computer science, and every sound rule
is a theorem in working clothes. Then a theorem would be found under every sound rule. The
census finds theorems under idealisations 6 and 4 and none under 7 and 8, which between
them generate requirements, operations, maintenance and migration. And of the four
practices examined, the practice that does rest on a theorem is the practice the theorem cannot
license. Walls are
theorems, and walls do not choose among the options they leave open.

Suppose instead that engineers are not minimising cost at all, but doing what social
organisms do: keeping their standing, following taste, protecting the story of who they
are. The paper takes this explanation most seriously, because much of it is true. Category (d)
of the test is that observation written down, and \S\ref{sec:whykeydropped} is about rules
held for social reasons after the cost reasons have gone. What it cannot explain is the direction
of change. The migration write-ups record engineers overruling their own
group's habit and giving a measured number as the reason.  Standing decides
which rules a person holds between crossings; it does not decide where the crossings are.
And taste, examined, is a rule whose conditions the holder can no longer state, which is a
fact about memory rather than about the absence of cost.

Suppose last that the machine does not matter, and software is social all the way down.
Then nothing about the machine would show up in every system regardless of who built it.
But every machine larger than its fastest memory has a hierarchy, and the reason it cannot
be engineered away is a theorem about distance. Every system that spans a network must
choose what it gives up during a partition, by a theorem about messages. And the three
bills that bite, moving, failing and stepping, are paid by every team on every continent.
The machine sets the floor and the people build on it, and an explanation that dropped
either half would explain half the field.

% ===========================================================================
\section{What This Forbids}\label{sec:forbids}
% ===========================================================================
An account that forbids nothing is worth little. Each of the following is a sentence
somebody currently says, and the account rules it out. ``This is a best practice'', with no
answer to Test~\ref{test}, is an assertion of optimality with nothing named to be optimal
for. ``This practice is universal'' means one of the three things the record's inversion field
lists, and the holder should say which. ``We tried it
and it worked'' is no evidence about the boundary, because under a favourable regime
nearly anything works. ``Physics requires it'' is ruled out above the four bills; multicore
is the standing example. ``This is a judgment call'' as a terminal explanation is ruled
out, because every judgment call has an operation, a row and an idealisation set, and
naming them is where the explanation begins. Verifying machine-written code by reading it
is ruled out as a default, because it verifies the code, whose cost went to zero. Treating
dependency management as a lookup is ruled out, because the failures come from 5 and 8 and
the practice addresses 6.

% ===========================================================================
\section{Threats to Validity}\label{sec:limits}
% ===========================================================================
The completeness of the eight, and so of the grammar, is a failed search, the position
Turing occupied with respect to his own thesis. The largest weakness is that the census of \S\ref{sec:census} has not been run: the
paper's central empirical claim could have been tested before publication, in days, and
was not. The residue coding of \S\ref{sec:verification} is a coding of a catalogue, not a
measurement. The latency bound is verified from the two published abstracts, which state
the physical assumptions and call the hierarchy's latency optimal, not from the full
proof. The inversions table is a curated sample, not a population, and two counter-cases
are recorded. The four verdicts are four; a fifth could go the other way. No single decision is followed
through the grammar, the test, the formula and the record on one real system with real
numbers, which is the largest gap in explanation as the unrun census is the largest gap in
evidence. The term software
intention overlaps the reference model's specification and could be argued to be nothing
more; if so, the claim that survives is that the software intention has the coordinates given
here. 

% ===========================================================================
\section{Predictions}\label{sec:predictions}
% ===========================================================================
The hypotheses of
\S\ref{sec:forward} carry their own observables; these are the predictions the framework
makes about itself.

\begin{pred}[Every decision codes onto the eight]
A sample of decisions from a project's history, coded by operation, row and idealisation
set, will code completely, and the hardest decisions will have the largest sets. A
substantial uncodable residue refutes the grammar.
\end{pred}

\begin{pred}[Every two-way practice has an inequality]
For any practice with documented two-way traffic, two nameable, measurable costs govern
it, and the inversion follows the crossover computed from independent cost series by a lag
of years; Sutter's self-reported lag of one to two years is one data point
\cite{sutter2005}. Failure to find the two costs on a majority of the cases of
\S\ref{sec:inversions} refutes the central claim.
\end{pred}

\begin{pred}[Genre predicts whether conditions survive]
Practices published with a formula get recomputed; practices published as imperatives get
argued about; the correlation is with the form of the original document, not the subject.
\end{pred}

\begin{pred}[The condition gap is large and socially sourced]
Practitioners asked to state a practice they follow will do so immediately; asked when it
would be wrong they will do markedly worse, and the gap will be larger for practices
learned socially than for practices derived themselves.
\end{pred}

\begin{pred}[Halt cost dominates change cost]
The cost of a change is dominated not by the size of the diff but by the halt cost of the
highest row of \S\ref{sec:rungs} at which it binds a decision.
\end{pred}

\begin{pred}[Records beat imperatives, for machines]
An agent given practice records with ratios and observables will outperform an agent given
the same practices as imperatives, with the margin concentrated where the regime does not
hold. A margin where the regime does hold means the mechanism proposed here is wrong.
\end{pred}

\begin{pred}[The residue moves in one direction]
When a machine doing engineering work is deployed and every human intervention is coded by
idealisation, the share at 1 to 6 falls as tooling improves and the share at 7 and 8 does
not. If it also falls, preferences and forecasts are recoverable from artefacts, and
\S\ref{sec:stops} is wrong.
\end{pred}

Take any principle software engineering teaches and try to name what it is answerable to under
Test~\ref{test} and which idealisations its alternatives pay under. Most will code. Some may not: a principle for which, after looking, no wall, no bill, no
idealisation, no measurement and no convention can be found. Those are the only things in
software engineering this paper cannot explain. Either they reveal a ninth cost, or they reveal the
interpreter doing something other than minimising cost.

% ===========================================================================
\section{Conclusion}
% ===========================================================================
Turing described a computer exactly by setting every cost of computing to zero, and
software engineering grew in the space he cleared. The engineer is an interpreter between people who
want things and a machine that does what its sheet says, and the software intention the
interpreter produces is shaped by a few walls and by the prices of everything else, the
eight costs the model idealised away. Nothing is chosen in a machine where nothing costs anything, so every real decision pays
differently under those eight, and the principles of software engineering are the cheaper sides
of inequalities whose conditions can be recovered from the inventory. A decision costs
writing, checking, the chance of being wrong times the cost of undoing, and finding out.
The cost falling fastest is the first, the one physics never priced; when it goes, what is
left is about being wrong, the engineer's place is where the missing input is a preference
or a forecast, and what gets verified changes from the code to the decision. The claims are held as Turing held his thesis, by a search for exceptions, and the
exceptions are the invitation.

% ===========================================================================
\section*{Supplementary Materials}
% ===========================================================================
A companion document, \emph{How the Sources Were Read}, records for every reference below
where in this paper it is used, what it was taken to support, how it was interpreted, and
where that interpretation could be mistaken; where the paper leaned on a formal result
without explaining it, the explanation is there. It is published with the paper as a PDF
and as a page where each entry can be opened, the source followed, and an objection filed
against that entry alone; the same page lets a reader bring a situation and have the
paper's concepts applied to it, and file what did not fit. Objections and exceptions that
survive change the paper, with credit. Paper and companion live at
\url{https://github.com/Onyangoe21/conceptual-space-of-software-architecture} and are
rendered at \url{https://curiousedwin.com/publications/what-computation-costs}.

% ===========================================================================
\section*{References}
% ===========================================================================
\footnotesize
\begin{thebibliography}{999}
\sloppy
\emergencystretch=3em
\setlength{\itemsep}{1.5pt}
\bibitem{abran2006} Abran, A., Meridji, K., Dolado, J.J. (2006). Software engineering from an engineering perspective: SWEBOK as a study object. \emph{UPGRADE/Nov\'atica} 7(1):46--52.
\bibitem{acton2014} Acton, M. (2014). Data-oriented design and C\texttt{++}. CppCon 2014. Slide text from the official deck.
\bibitem{aggarwal1987} Aggarwal, A., Alpern, B., Chandra, A.K., Snir, M. (1987). A model for hierarchical memory. \emph{Proc.\ STOC 1987}.
\bibitem{alexander1977} Alexander, C., Ishikawa, S., Silverstein, M. (1977). \emph{A Pattern Language}. Oxford University Press.
\bibitem{amdahl1967} Amdahl, G.M. (1967). Validity of the single processor approach to achieving large scale computing capabilities. \emph{AFIPS Conf.\ Proc.} 30:483--485.
\bibitem{anderson1982} Anderson, J.R. (1982). Acquisition of cognitive skill. \emph{Psychological Review} 89(4):369--406.
\bibitem{anderson2016} Anderson, D.J., Carmichael, A. (2016). \emph{Essential Kanban Condensed}. Lean Kanban University Press.
\bibitem{appuswamy2019} Appuswamy, R., Borovica-Gajic, R., Graefe, G., Ailamaki, A. (2019). The five-minute rule 30 years later. \emph{Communications of the ACM} 62(11).
\bibitem{arthur1989} Arthur, W.B. (1989). Competing technologies, increasing returns, and lock-in by historical events. \emph{Economic Journal} 99:116--131.
\bibitem{bennett1973} Bennett, C.H. (1973). Logical reversibility of computation. \emph{IBM Journal of Research and Development} 17(6):525--532.
\bibitem{bernard2002} Bernard, M., Fernandez, M., Hull, S. (2002). The effects of line length on children and adults' online reading performance. \emph{Usability News} 4(2).
\bibitem{berut2012} B\'erut, A., et al. (2012). Experimental verification of Landauer's principle linking information and thermodynamics. \emph{Nature} 483:187--189.
\bibitem{bilardi2009} Bilardi, G., Ekanadham, K., Pattnaik, P. (2009). On approximating the ideal random access machine by physical machines. \emph{Journal of the ACM} 56(5), Article 27. Earlier: \emph{Proc.\ SPAA 2002}, pp.~109--116, which states the $\Omega(x^{1/d})$ lower bound.
\bibitem{bozkurt2017} Bozkurt, I.N., et al. (2017). Why is the Internet so slow?! \emph{Proc.\ PAM 2017}.
\bibitem{brooks1987} Brooks, F.P. (1987). No silver bullet: essence and accidents of software engineering. \emph{IEEE Computer} 20(4):10--19.
\bibitem{buse2010} Buse, R.P.L., Weimer, W. (2010). Learning a metric for code readability. \emph{IEEE TSE} 36(4):546--558.
\bibitem{c11} ISO/IEC 9899:2011, draft N1570. Observable behaviour \S5.1.2.3~\P6.
\bibitem{chipsandcheese} Lam, C. (2022). AMD's Zen 4, part 2: memory subsystem. Chips and Cheese, 8 November 2022.
\bibitem{church1936} Church, A. (1936). An unsolvable problem of elementary number theory. \emph{American Journal of Mathematics} 58(2):345--363.
\bibitem{codd1970} Codd, E.F. (1970). A relational model of data for large shared data banks. \emph{Communications of the ACM} 13(6):377--387.
\bibitem{conway1968} Conway, M.E. (1968). How do committees invent? \emph{Datamation} 14(4):28--31.
\bibitem{dally2022} Dally, W. (2022). On the model of computation: point. \emph{Communications of the ACM} 65(9):30--32.
\bibitem{david1985} David, P.A. (1985). Clio and the economics of QWERTY. \emph{American Economic Review} 75(2):332--337.
\bibitem{dennard1974} Dennard, R.H., et al. (1974). Design of ion-implanted MOSFETs with very small physical dimensions. \emph{IEEE JSSC} 9(5):256--268.
\bibitem{devanbu2016} Devanbu, P., Zimmermann, T., Bird, C. (2016). Belief \& evidence in empirical software engineering. \emph{Proc.\ ICSE 2016}.
\bibitem{dhh2023} Heinemeier Hansson, D. (2023). We have left the cloud. world.hey.com, 23 June 2023.
\bibitem{dyson2004} Dyson, M.C. (2004). How physical text layout affects reading from screen. \emph{Behaviour \& Information Technology} 23(6):377--393.
\bibitem{earman1999} Earman, J., Norton, J.D. (1998--99). Exorcist XIV: the wrath of Maxwell's demon. \emph{Studies in History and Philosophy of Modern Physics}, parts I and II.
\bibitem{eswaran1976} Eswaran, K.P., Gray, J.N., Lorie, R.A., Traiger, I.L. (1976). The notions of consistency and predicate locks in a database system. \emph{Communications of the ACM} 19(11):624--633.
\bibitem{fallacies} Deutsch, L.P., and colleagues (c.~1994). The fallacies of distributed computing. Explicated in Rotem-Gal-Oz, A. (2006).
\bibitem{flp1985} Fischer, M.J., Lynch, N.A., Paterson, M.S. (1985). Impossibility of distributed consensus with one faulty process. \emph{Journal of the ACM} 32(2):374--382.
\bibitem{gamma1994} Gamma, E., Helm, R., Johnson, R., Vlissides, J. (1994). \emph{Design Patterns}. Addison-Wesley.
\bibitem{gandy1980} Gandy, R. (1980). Church's thesis and principles for mechanisms. In \emph{The Kleene Symposium}. North-Holland.
\bibitem{gilbertlynch2002} Gilbert, S., Lynch, N. (2002). Brewer's conjecture and the feasibility of consistent, available, partition-tolerant web services. \emph{ACM SIGACT News} 33(2):51--59.
\bibitem{gluck2020} Gluck, A. (2020). Introducing domain-oriented microservice architecture. Uber Engineering Blog, 23 July 2020.
\bibitem{graefe2007} Graefe, G. (2007). The five-minute rule twenty years later. \emph{Proc.\ DaMoN 2007}; republished \emph{ACM Queue} 6(4), 2008.
\bibitem{gray1987} Gray, J., Putzolu, G.R. (1987). The 5 minute rule for trading memory for disc accesses. \emph{Proc.\ ACM SIGMOD 1987}, pp.~395--398.
\bibitem{graygraefe1997} Gray, J., Graefe, G. (1997). The five-minute rule ten years later. \emph{ACM SIGMOD Record} 26(4):63--68.
\bibitem{grier2005} Grier, D.A. (2005). \emph{When Computers Were Human}. Princeton University Press.
\bibitem{grudin1996} Grudin, J. (1996). Evaluating opportunities for design capture. In \emph{Design Rationale}. Lawrence Erlbaum.
\bibitem{gunter2000} Gunter, C.A., Gunter, E.L., Jackson, M., Zave, P. (2000). A reference model for requirements and specifications. \emph{IEEE Software} 17(3):37--43.
\bibitem{gustafson1988} Gustafson, J.L. (1988). Reevaluating Amdahl's law. \emph{Communications of the ACM} 31(5):532--533.
\bibitem{hayden2025} Hayden, P., Wang, J. (2025). On the Bekenstein bound. \emph{Quantum} 9:1664.
\bibitem{hennessy2019} Hennessy, J.L., Patterson, D.A. (2019). A new golden age for computer architecture. \emph{Communications of the ACM} 62(2):48--60.
\bibitem{hodges1983} Hodges, A. (1983). \emph{Alan Turing: The Enigma}. Burnett Books.
\bibitem{hong2016} Hong, J., Lambson, B., Dhuey, S., Bokor, J. (2016). Experimental test of Landauer's principle in single-bit operations on nanomagnetic memory bits. \emph{Science Advances} 2(3):e1501492.
\bibitem{hpqa} Hennessy, J.L., Patterson, D.A. \emph{Computer Architecture: A Quantitative Approach}, ch.~2.
\bibitem{ibm-card} IBM. The punched card. ibm.com/history/punched-card.
\bibitem{jansenbosch2005} Jansen, A., Bosch, J. (2005). Software architecture as a set of architectural design decisions. \emph{Proc.\ WICSA 2005}.
\bibitem{jones-cards} Jones, D.W. Punched cards: a brief illustrated technical history. University of Iowa.
\bibitem{jun2014} Jun, Y., Gavrilov, M., Bechhoefer, J. (2014). High-precision test of Landauer's principle in a feedback trap. \emph{Physical Review Letters} 113:190601.
\bibitem{kalia2019} Kalia, A., Kaminsky, M., Andersen, D.G. (2019). Datacenter RPCs can be general and fast. \emph{Proc.\ NSDI 2019}.
\bibitem{kingman1961} Kingman, J.F.C. (1961). The single server queue in heavy traffic. \emph{Math.\ Proc.\ Cambridge Phil.\ Soc.} 57(4):902--904.
\bibitem{kleene1936} Kleene, S.C. (1936). Lambda-definability and recursiveness. \emph{Duke Mathematical Journal} 2:340--353.
\bibitem{klein1986} Klein, G.A., Calderwood, R., Clinton-Cirocco, A. (1986). Rapid decision making on the fire ground. \emph{Proc.\ Human Factors Society} 30(16):576--580.
\bibitem{knuth1974} Knuth, D.E. (1974). Structured programming with go to statements. \emph{ACM Computing Surveys} 6(4):261--301. The passage is at p.~268.
\bibitem{landauer1961} Landauer, R. (1961). Irreversibility and heat generation in the computing process. \emph{IBM Journal of Research and Development} 5(3):183--191.
\bibitem{lane2023} Lane, T. (2023). Re: let's make PostgreSQL multi-threaded. pgsql-hackers, 5 June 2023.
\bibitem{lehman1980} Lehman, M.M. (1980). Programs, life cycles, and laws of software evolution. \emph{Proceedings of the IEEE} 68(9):1060--1076.
\bibitem{leiserson2020} Leiserson, C.E., et al. (2020). There's plenty of room at the Top. \emph{Science} 368(6495):eaam9744.
\bibitem{li2007} Li, C., Ding, C., Shen, K. (2007). Quantifying the cost of context switch. \emph{Proc.\ ExpCS 2007}.
\bibitem{little1961} Little, J.D.C. (1961). A proof for the queuing formula: $L=\lambda W$. \emph{Operations Research} 9(3):383--387.
\bibitem{lloyd2000} Lloyd, S. (2000). Ultimate physical limits to computation. \emph{Nature} 406:1047--1054.
\bibitem{margolus1998} Margolus, N., Levitin, L.B. (1998). The maximum speed of dynamical evolution. \emph{Physica D} 120:188--195.
\bibitem{mcmullin1985} McMullin, E. (1985). Galilean idealization. \emph{Studies in History and Philosophy of Science} 16(3):247--273.
\bibitem{mendonca2021} Mendon\c{c}a, N.C., Box, C., Manolache, C., Ryan, L. (2021). The monolith strikes back. \emph{IEEE Software} 38:17--22.
\bibitem{menzies2017} Menzies, T., Nichols, W., Shull, F., Layman, L. (2017). Are delayed issues harder to resolve? \emph{Empirical Software Engineering} 22(4):1903--1935.
\bibitem{meza2015} Meza, J., Wu, Q., Kumar, S., Mutlu, O. (2015). A large-scale study of flash memory failures in the field. \emph{Proc.\ SIGMETRICS 2015}.
\bibitem{nanavati2005} Nanavati, A.A., Bias, R.G. (2005). Optimal line length in reading: a literature review. \emph{Visible Language} 39(2):121--145.
\bibitem{napkinmath} Eskildsen, S. \emph{napkin-math}. Single-host rows re-measured on GCP \texttt{c4-standard-48-lssd}, 8 March 2026.
\bibitem{naur1985} Naur, P. (1985). Programming as theory building. \emph{Microprocessing and Microprogramming} 15(5):253--261.
\bibitem{noonan2018} Noonan, A. (2018). Goodbye microservices. Segment Engineering Blog, 10 July 2018.
\bibitem{norton2005} Norton, J.D. (2005). Eaters of the lotus: Landauer's principle and the return of Maxwell's demon. \emph{Studies in History and Philosophy of Modern Physics} 36:375--411.
\bibitem{nygard2011} Nygard, M. (2011). Documenting architecture decisions. cognitect.com.
\bibitem{parnas1972} Parnas, D.L. (1972). On the criteria to be used in decomposing systems into modules. \emph{Communications of the ACM} 15(12):1053--1058.
\bibitem{prechelt2001} Prechelt, L., Unger, B., Tichy, W.F., Br\"ossler, P., Votta, L.G. (2001). A controlled experiment in maintenance comparing design patterns to simpler solutions. \emph{IEEE TSE} 27(12):1134--1144.
\bibitem{proesmans2020} Proesmans, K., Ehrich, J., Bechhoefer, J. (2020). Finite-time Landauer principle. \emph{Physical Review Letters} 125:100602.
\bibitem{reis2019} Reis, C., Moshchuk, A., Oskov, N. (2019). Site isolation. \emph{Proc.\ USENIX Security 2019}, pp.~1661--1678.
\bibitem{rice1953} Rice, H.G. (1953). Classes of recursively enumerable sets and their decision problems. \emph{Transactions of the American Mathematical Society} 74(2):358--366.
\bibitem{saltzer1984} Saltzer, J.H., Reed, D.P., Clark, D.D. (1984). End-to-end arguments in system design. \emph{ACM TOCS} 2(4):277--288.
\bibitem{schroeder2009} Schroeder, B., Pinheiro, E., Weber, W.-D. (2009). DRAM errors in the wild: a large-scale field study. \emph{Proc.\ SIGMETRICS 2009}.
\bibitem{sevencpu} 7-cpu.com, Intel Skylake page. i7-6700 at 4.0\,GHz, DDR4-2400 CL15.
\bibitem{shipman1999} Shipman, F.M., Marshall, C.C. (1999). Formality considered harmful. \emph{Computer Supported Cooperative Work} 8:333--352.
\bibitem{shirriff2019} Shirriff, K. (2019). IBM, sonic delay lines, and the history of the 80$\times$24 display. righto.com, November 2019.
\bibitem{shute2013} Shute, J., et al. (2013). F1: a distributed SQL database that scales. \emph{PVLDB} 6(11):1068--1079.
\bibitem{sieg1997} Sieg, W. (1997). Step by recursive step: Church's analysis of effective calculability. \emph{Bulletin of Symbolic Logic} 3(2):154--180.
\bibitem{simon1962} Simon, H.A. (1962). The architecture of complexity. \emph{Proc.\ American Philosophical Society} 106(6):467--482.
\bibitem{simon1996} Simon, H.A. (1996). \emph{The Sciences of the Artificial}, 3rd ed. MIT Press.
\bibitem{sjoberg2018} Sj\o berg, D.I.K. (2018). An empirical study of WIP in kanban teams. \emph{Proc.\ ESEM 2018}.
\bibitem{soare1996} Soare, R.I. (1996). Computability and recursion. \emph{Bulletin of Symbolic Logic} 2(3):284--321.
\bibitem{soares2010} Soares, L., Stumm, M. (2010). FlexSC: flexible system call scheduling with exception-less system calls. \emph{Proc.\ OSDI 2010}.
\bibitem{stidham1974} Stidham, S., Jr. (1974). A last word on $L=\lambda W$. \emph{Operations Research} 22(2):417--421.
\bibitem{stonebraker2005} Stonebraker, M., \c{C}etintemel, U. (2005). ``One size fits all'': an idea whose time has come and gone. \emph{Proc.\ ICDE 2005}.
\bibitem{sutter2005} Sutter, H. (2005). The free lunch is over. \emph{Dr.\ Dobb's Journal} 30(3).
\bibitem{turing1936} Turing, A.M. (1936--37). On computable numbers, with an application to the Entscheidungsproblem. \emph{Proc.\ LMS} s2-42:230--265.
\bibitem{turing1950} Turing, A.M. (1950). Computing machinery and intelligence. \emph{Mind} 59(236):433--460.
\bibitem{vacanti2015} Vacanti, D. (2015). \emph{Actionable Agile Metrics for Predictability}. ActionableAgile Press.
\bibitem{valiant1990} Valiant, L.G. (1990). A bridging model for parallel computation. \emph{Communications of the ACM} 33(8):103--111.
\bibitem{vanleeuwen2001} van Leeuwen, J., Wiedermann, J. (2001). The Turing machine paradigm in contemporary computing. In \emph{Mathematics Unlimited: 2001 and Beyond}. Springer.
\bibitem{vincenti1990} Vincenti, W.G. (1990). \emph{What Engineers Know and How They Know It}. Johns Hopkins University Press.
\bibitem{vonneumann1956} von Neumann, J. (1956). Probabilistic logics and the synthesis of reliable organisms from unreliable components. In \emph{Automata Studies}. Princeton University Press.
\bibitem{wulf1995} Wulf, W.A., McKee, S.A. (1995). Hitting the memory wall: implications of the obvious. \emph{ACM SIGARCH CAN} 23(1):20--24.
\bibitem{yan2018} Yan, L.L., et al. (2018). Single-atom demonstration of the quantum Landauer principle. \emph{Physical Review Letters} 120:210601.
\bibitem{zannier2007} Zannier, C., Chiasson, M., Maurer, F. (2007). A model of design decision making based on empirical results of interviews with software designers. \emph{Information and Software Technology} 49(6):637--653.
\bibitem{zavejackson1997} Zave, P., Jackson, M. (1997). Four dark corners of requirements engineering. \emph{ACM Transactions on Software Engineering and Methodology} 6(1):1--30.
\bibitem{zhang2025} Zhang, T., et al. (2025). Five-minute rule 40 years later. arXiv:2511.03944.
\end{thebibliography}

\end{document}
